\chapter{Eine Veränderung von außen}

\section{Der Gedanke}

\begin{merksatz}
\boxtitel{Was man braucht}
Die Zuzahlung ist endogen. Gesucht ist eine Größe $z$, die sie verschiebt, ohne selbst mit der Gesundheit zusammenzuhängen.

\begin{center}
{\small
\renewcommand{\arraystretch}{1.55}
\begin{tabular}{@{}P{2.8cm}P{5.8cm}@{}}
\toprule
\textbf{Relevanz} & $\operatorname{Cov}(z,x) \ne 0$ --- prüfbar.\\
\textbf{Exogenität} & $\operatorname{Cov}(z,u) = 0$ --- nicht prüfbar.\\
\bottomrule
\end{tabular}}
\end{center}

\begin{center}
\begin{minipage}{0.88\textwidth}
\kursiv{\textbf{Das ist das Instrument aus Band 23, Kapitel~9} --- dort ein Ausweg, wenn keine Anpassung gelingt, hier der Normalfall einer ganzen Disziplin. Der Unterschied liegt in der Herkunft: Die Epidemiologie sucht Instrumente in der Natur, die Ökonometrie findet sie meist in Gesetzen, Stichtagen und Verwaltungsregeln.}
\end{minipage}
\end{center}

\schritt{Im Leitfall}
Eine Reform erhöht die Zuzahlung in einem Teil der Kassen, aus Gründen, die mit dem Gesundheitszustand der Versicherten nichts zu tun haben.
\end{merksatz}

\begin{merksatz}
\boxtitel{Satz 3.1 --- der Instrumentschätzer}
Erfüllt $z$ beide Bedingungen, so gilt im Modell $y = \beta x + u$
\[ \hat\beta_{\text{IV}} = \frac{\operatorname{Cov}(z,y)}{\operatorname{Cov}(z,x)} \ \longrightarrow\ \beta . \]

\schritt{Beweis}
\[ \operatorname{Cov}(z,y) = \operatorname{Cov}(z,\beta x + u) = \beta\operatorname{Cov}(z,x) + \underbrace{\operatorname{Cov}(z,u)}_{=\,0} \]
Division durch $\operatorname{Cov}(z,x) \ne 0$ liefert $\beta$. $\blacksquare$

\begin{center}
\begin{minipage}{0.9\textwidth}
\textbf{Man benutzt nur denjenigen Teil der Schwankung von $x$, der von $z$ stammt --- und der ist nach Annahme unschuldig.}
\end{minipage}
\end{center}
\end{merksatz}

\section{Zweistufige Schätzung}

\begin{merksatz}
\boxtitel{Dasselbe in zwei Regressionen}
\begin{center}
{\small
\renewcommand{\arraystretch}{1.55}
\begin{tabular}{@{}P{2.4cm}P{6.2cm}@{}}
\toprule
\textbf{Erste Stufe} & $x$ auf $z$ (und alle exogenen Größen) regressieren; die angepassten Werte $\hat x$ behalten.\\
\textbf{Zweite Stufe} & $y$ auf $\hat x$ regressieren.\\
\bottomrule
\end{tabular}}
\end{center}

\begin{center}
\begin{minipage}{0.88\textwidth}
\kursiv{\textbf{Der Vorzug gegenüber Satz 3.1: Das Verfahren verträgt mehrere Instrumente und mehrere Regressoren.} Bei einem Instrument und einem Regressor ergibt es genau denselben Wert --- $\hat x$ ist dann eine lineare Funktion von $z$, und die Kovarianzen kürzen sich.}
\end{minipage}
\end{center}

\schritt{Im Leitfall gerechnet}
$200\,000$ Versicherte, die Reform erhöht die Zuzahlung um $7{,}83$ Euro:

\begin{center}
{\small
\renewcommand{\arraystretch}{1.5}
\begin{tabular}{@{}P{3.4cm}P{2.4cm}@{}}
\toprule
\textbf{Verfahren} & \textbf{je Euro}\\
\midrule
kleinste Quadrate & $-0{,}297$\\
zweistufig & $\mathbf{-0{,}100}$\\
Quotient nach Satz 3.1 & $-0{,}100$\\
\midrule
wahrer Wert & $\mathbf{-0{,}100}$\\
\bottomrule
\end{tabular}}
\end{center}

\begin{center}
\begin{minipage}{0.9\textwidth}
\textbf{Eine einzige exogene Veränderung genügt --- die Gesundheit muss gar nicht gemessen werden.}
\end{minipage}
\end{center}

\begin{center}
\begin{minipage}{0.88\textwidth}
\kursiv{\textbf{Das ist der Grund für die Anziehungskraft des Verfahrens.} Es verlangt nichts über die Störgröße --- weder ihre Messung noch ihre Form --- sondern nur, dass eine bestimmte andere Größe mit ihr nichts zu tun hat. Man tauscht eine unprüfbare Annahme gegen eine andere, und die zweite ist oft leichter zu verteidigen.}
\end{minipage}
\end{center}
\end{merksatz}

\section{Der Bruchfall}

\begin{warnung}
\boxtitelr{Schwache Instrumente}
Dieselbe Rechnung mit $2000$ Versicherten und schwächer werdender Reform, je $500$ Wiederholungen:

\begin{center}
{\small
\renewcommand{\arraystretch}{1.48}
\begin{tabular}{@{}P{2.0cm}P{1.4cm}P{1.8cm}P{3.0cm}@{}}
\toprule
\textbf{Wirkung} & $F$ & \textbf{Median} & \textbf{$5$--$95$ Prozent}\\
\midrule
stark & $719$ & $-0{,}098$ & $[-0{,}121;\ -0{,}073]$\\
mäßig & $44$ & $-0{,}100$ & $[-0{,}177;\ +0{,}019]$\\
schwach & $2{,}8$ & $-0{,}135$ & $\mathbf{[-0{,}805;\ +0{,}922]}$\\
\bottomrule
\end{tabular}}
\end{center}
\kursiv{Zum Vergleich: Die kleinsten Quadrate liefern stets etwa $-0{,}30$.}

\begin{center}
\begin{minipage}{0.9\textwidth}
\textbf{Bei schwachem Instrument reicht das Streuungsband von $-0{,}8$ bis $+0{,}9$ --- die Schätzung ist unbrauchbar, ohne dass es ihr anzusehen wäre.}
\end{minipage}
\end{center}

\schritt{Die beiden Gefahren}
\begin{center}
{\small
\renewcommand{\arraystretch}{1.52}
\begin{tabular}{@{}P{2.6cm}P{6.0cm}@{}}
\toprule
\textbf{Streuung} & Der Nenner ist klein und selbst zufällig --- der Quotient schwankt wild.\\
\textbf{Verzerrung} & Bei schwachem Instrument wandert die Schätzung zum Wert der kleinsten Quadrate hin --- genau dorthin, wovon man wegwollte.\\
\bottomrule
\end{tabular}}
\end{center}

\begin{center}
\begin{minipage}{0.88\textwidth}
\kursiv{\textbf{Die Faustregel $F > 10$ stammt aus einer Arbeit von 1997 und gilt heute als zu großzügig}; neuere Untersuchungen verlangen je nach Lage Werte über hundert. Entscheidend ist, dass $F$ \kursiv{vor} der Schätzung berichtet wird --- und dass ein kleines $F$ nicht zu mehr Instrumenten verleitet, denn jedes schwache zusätzliche verschlimmert die Verzerrung.}
\end{minipage}
\end{center}
\end{warnung}

\begin{warnung}
\boxtitelr{Und die Bedingung, die niemand prüft}
\begin{center}
\begin{minipage}{0.88\textwidth}
\kursiv{Die Exogenität des Instruments ist nicht prüfbar --- außer man hat mehr Instrumente als nötig. \textbf{Dann vergleicht man die Schätzungen}: Weichen sie stärker voneinander ab, als der Zufall erlaubt, so ist mindestens eines ungültig. Das ist der Überidentifikationstest, und er sagt nicht, welches.}
\end{minipage}
\end{center}

\schritt{Gerechnet}
Zwei Instrumente, von denen das zweite unerlaubt unmittelbar auf die Besuche wirkt:

\begin{center}
{\small
\renewcommand{\arraystretch}{1.48}
\begin{tabular}{@{}P{2.6cm}P{2.4cm}P{2.4cm}@{}}
\toprule
\textbf{Instrument} & \textbf{Schätzung} & \\
\midrule
$z_1$ (gültig) & $-0{,}098$ & \textcolor{gruen}{richtig}\\
$z_2$ (ungültig) & $-0{,}013$ & \textcolor{akzent}{falsch}\\
\midrule
Differenz & $\mathbf{0{,}085}$ & Warnsignal\\
\bottomrule
\end{tabular}}
\end{center}

\begin{center}
\begin{minipage}{0.88\textwidth}
\kursiv{\textbf{Besteht man den Test, so ist nichts bewiesen} --- alle Instrumente können gemeinsam in dieselbe Richtung verletzt sein, und dann stimmen die Schätzungen überein, ohne richtig zu sein. Der Test ist eine notwendige Bedingung, wie der Balancetest aus Band 23, Kapitel~7 und der Dichtetest aus Kapitel~10 desselben Bandes.}
\end{minipage}
\end{center}
\end{warnung}

\section{Übungsaufgaben}

\begin{aufgabe}
\boxtitel{Aufgabe 3.1 --- Den Schätzer von Hand}
Aus einer Reform: In Kassen ohne Reform zahlen die Versicherten im Mittel $10{,}4$ Euro und gehen $6{,}95$-mal zum Arzt; in Kassen mit Reform $18{,}2$ Euro und $6{,}17$-mal.

\loesung
\schritt{Schritt 1 --- Die beiden Differenzen}
\[ \Delta x = 18{,}2 - 10{,}4 = 7{,}8 \ \text{Euro}, \qquad \Delta y = 6{,}17 - 6{,}95 = -0{,}78 \ \text{Besuche} \]

\schritt{Schritt 2 --- Der Quotient}
\[ \hat\beta = \frac{-0{,}78}{7{,}8} = \mathbf{-0{,}100} \]

\schritt{Schritt 3 --- Die Deutung der beiden Teile}
\begin{center}
{\small
\renewcommand{\arraystretch}{1.5}
\begin{tabular}{@{}P{2.6cm}P{6.0cm}@{}}
\toprule
\textbf{Zähler} & Wirkung der Reform auf die Besuche --- die Absicht, wie in Band 23, Kapitel~3.\\
\textbf{Nenner} & Wirkung der Reform auf die Zuzahlung --- der Befolgungsgrad.\\
\bottomrule
\end{tabular}}
\end{center}

\begin{center}
\begin{minipage}{0.88\textwidth}
\kursiv{\textbf{Der Schätzer ist die Wirkung der Anordnung, geteilt durch ihre Durchsetzung} --- derselbe Ausdruck wie der Wald-Schätzer aus Band 23. Dort hieß der Nenner Befolgeranteil, hier erste Stufe; die Rechnung ist dieselbe.}
\end{minipage}
\end{center}

\schritt{Schritt 4 --- Die Zielgruppe}
\begin{center}
\begin{minipage}{0.88\textwidth}
\kursiv{Wie dort gilt das Ergebnis für die, deren Zuzahlung sich durch die Reform geändert hat. \textbf{Wer ohnehin einen Hochzuzahlungstarif hatte, trägt nichts bei} --- und wenn solche Versicherten anders reagieren, gilt die Zahl für sie nicht.}
\end{minipage}
\end{center}
\end{aufgabe}

\begin{aufgabe}
\boxtitel{Aufgabe 3.2 --- Warum die zweite Stufe falsche Standardfehler liefert}
Wer die zweite Stufe als gewöhnliche Regression rechnet, erhält zu kleine Streuungen. Warum?

\loesung
\schritt{Schritt 1 --- Der Fehler}
Die gewöhnliche Regression auf $\hat x$ benutzt die Residuen $y - \hat\beta\hat x$. Richtig wären die Residuen $y - \hat\beta x$ --- mit dem \kursiv{tatsächlichen} $x$.

\schritt{Schritt 2 --- Der Unterschied}
\begin{center}
{\small
\renewcommand{\arraystretch}{1.5}
\begin{tabular}{@{}P{2.8cm}P{5.8cm}@{}}
\toprule
\textbf{Falsch} & Die Streuung des Modells wird an den geglätteten Werten gemessen und erscheint kleiner.\\
\textbf{Richtig} & Jedes Programm für zweistufige Schätzung rechnet das automatisch.\\
\bottomrule
\end{tabular}}
\end{center}

\begin{center}
\begin{minipage}{0.9\textwidth}
\textbf{Die beiden Stufen von Hand zu rechnen, liefert den richtigen Koeffizienten und einen falschen Standardfehler.}
\end{minipage}
\end{center}

\schritt{Schritt 3 --- Und die Größenordnung}
\begin{center}
\begin{minipage}{0.88\textwidth}
\kursiv{Der Fehler fällt umso stärker aus, je weniger von $x$ das Instrument erklärt --- also gerade bei schwachen Instrumenten, wo die Streuung ohnehin das Problem ist. \textbf{Beide Fehler wirken in dieselbe Richtung: Man unterschätzt die Unsicherheit dort am meisten, wo sie am größten ist.}}
\end{minipage}
\end{center}
\end{aufgabe}

\begin{aufgabe}
\boxtitel{Aufgabe 3.3 --- Instrumente beurteilen}
Welche der folgenden Größen taugt als Instrument für die Zuzahlung?
\begin{center}
{\small
\renewcommand{\arraystretch}{1.5}
\begin{tabular}{@{}P{1.0cm}P{7.4cm}@{}}
\toprule
(a) & Der Beitragssatz der Kasse.\\
(b) & Das Bundesland, in dem der Versicherte wohnt.\\
(c) & Ein Stichtag: Wer nach dem 1.\ Januar eintritt, erhält den neuen Tarif.\\
\bottomrule
\end{tabular}}
\end{center}

\loesung
\schritt{(a) Beitragssatz}
\begin{center}
{\small
\renewcommand{\arraystretch}{1.48}
\begin{tabular}{@{}P{2.4cm}P{5.8cm}@{}}
\toprule
Relevanz & gegeben\\
Exogenität & \textcolor{akzent}{verletzt} --- der Beitragssatz hängt von der Krankheitslast der Versicherten ab\\
\bottomrule
\end{tabular}}
\end{center}

\schritt{(b) Bundesland}
\begin{center}
{\small
\renewcommand{\arraystretch}{1.48}
\begin{tabular}{@{}P{2.4cm}P{5.8cm}@{}}
\toprule
Relevanz & gegeben, wenn die Tarife regional verschieden sind\\
Exogenität & \textcolor{akzent}{zweifelhaft} --- Bundesländer unterscheiden sich in Alter, Einkommen und Arztdichte\\
\bottomrule
\end{tabular}}
\end{center}

\schritt{(c) Stichtag}
\begin{center}
{\small
\renewcommand{\arraystretch}{1.48}
\begin{tabular}{@{}P{2.4cm}P{5.8cm}@{}}
\toprule
Relevanz & gegeben\\
Exogenität & \textcolor{gruen}{haltbar} --- wer am 28.\ Dezember eintritt, unterscheidet sich kaum von dem, der am 3.\ Januar eintritt\\
\bottomrule
\end{tabular}}
\end{center}

\begin{center}
\begin{minipage}{0.9\textwidth}
\textbf{Nur der Stichtag taugt --- und auch er nur in einem schmalen Fenster um den Termin.}
\end{minipage}
\end{center}

\schritt{Die Prüfung}
\begin{center}
\begin{minipage}{0.88\textwidth}
\kursiv{\textbf{Man prüft einen Stichtag wie eine Schwelle in Band 23, Kapitel~10}: Häufen sich die Eintritte unmittelbar vor dem Termin, so haben die Versicherten die Regel verstanden und sich sortiert --- dann ist die Exogenität dahin. \textbf{Dieselbe Rechnung, dasselbe Warnsignal, andere Fachsprache.} Die Ökonometrie nennt es Manipulation der Zuweisungsvariablen.}
\end{minipage}
\end{center}
\end{aufgabe}

\begin{geschichte}
\boxtitelg{Ein Anhang über Pflanzenöle}
\Person{Philip Wright} beschrieb den Instrumentschätzer 1928 im Anhang B eines Buches über Zölle auf Pflanzenöle --- Band 23, Kapitel~9 erzählt die Geschichte ausführlicher.

\begin{center}
\begin{minipage}{0.9\textwidth}
\kursiv{\textbf{Bemerkenswert ist, dass das Verfahren danach für Jahrzehnte verschwand}: Die Cowles-Kommission entwickelte in den 1940er Jahren die zweistufige Schätzung unabhängig und als Teil der Theorie simultaner Gleichungen. Erst in den 1990er Jahren wurde sie wieder als eigenständiges Werkzeug begriffen --- nicht als Lösung eines Systems, sondern als Ausnutzung einer einzelnen glaubwürdigen Veränderung.}
\end{minipage}
\end{center}

\schritt{Die Wende}
\begin{center}
\begin{minipage}{0.9\textwidth}
\kursiv{\Person{Joshua Angrist} und \Person{Alan Krueger} benutzten 1991 das Geburtsquartal als Instrument für die Schulbesuchsdauer --- weil die Schulpflicht an Geburtstage anknüpft. \textbf{Der Aufsatz machte Schule und wurde zugleich zum Lehrbeispiel für schwache Instrumente}: \Person{Bound}, \Person{Jaeger} und \Person{Baker} zeigten 1995, dass die Ergebnisse bei zufällig erzeugten Instrumenten fast genauso aussahen.}
\end{minipage}
\end{center}

\schritt{Und der Nobelpreis}
\begin{center}
\begin{minipage}{0.9\textwidth}
\kursiv{2021 erhielten \Person{Angrist} und \Person{Guido Imbens} den Preis für die Klärung, was Instrumente eigentlich schätzen --- den Effekt bei den Befolgern. \textbf{\Person{David Card} erhielt ihn zugleich für die Arbeiten zu natürlichen Versuchen}, und die Begründung nannte ausdrücklich die Verbindung von beidem: glaubwürdige Veränderung plus saubere Deutung dessen, was man misst.}
\end{minipage}
\end{center}
\end{geschichte}

\begin{forschung}
\boxtitelp{Instrumente heute}
\begin{center}
{\small
\renewcommand{\arraystretch}{1.45}
\begin{tabular}{@{}P{3.2cm}P{5.0cm}@{}}
\toprule
\textbf{Schwache Instrumente} & Verfahren, die auch bei kleinem $F$ gültige Intervalle liefern.\\
\textbf{Viele Instrumente} & Auswahl mit Lernverfahren --- und die Verzerrung, die daraus folgt.\\
\textbf{Fast gültig} & Schranken, wenn die Ausschlussbedingung nur ungefähr gilt.\\
\bottomrule
\end{tabular}}
\end{center}

\schritt{Zur ersten Zeile}
\begin{center}
\begin{minipage}{0.88\textwidth}
\kursiv{Es gibt Verfahren, deren Vertrauensintervalle unabhängig von der Instrumentstärke gültig bleiben --- der Anderson-Rubin-Test ist das bekannteste. \textbf{Bei sehr schwachen Instrumenten liefert er dann allerdings die ganze reelle Achse}, und genau das ist die richtige Antwort: Die Daten sagen nichts. Dass dieses Ergebnis selten berichtet wird, liegt nicht an der Methodik.}
\end{minipage}
\end{center}
\end{forschung}

\begin{merksatz}
\boxtitel{Was Sie jetzt können}
\begin{center}
\renewcommand{\arraystretch}{1.5}
\begin{tabular}{@{}P{3.0cm}P{5.6cm}@{}}
\toprule
\textbf{Satz 3.1} & $\operatorname{Cov}(z,y)/\operatorname{Cov}(z,x)$\\
\textbf{Zweistufig} & verträgt mehrere Instrumente\\
\textbf{Schwach} & Streuung und Verzerrung zugleich\\
\textbf{Überidentifiz.} & notwendige Bedingung, kein Beweis\\
\textbf{Stichtage} & die glaubwürdigsten Instrumente\\
\bottomrule
\end{tabular}
\end{center}

\begin{center}
\begin{minipage}{0.88\textwidth}
\kursiv{Damit endet Teil~I. Der Instrumentschätzer ist ein Sonderfall eines allgemeineren Gedankens: Man schreibt auf, welche Momente null sein müssen, und sucht die Parameter, die das erfüllen. Teil~II beginnt mit dieser Verallgemeinerung.}
\end{minipage}
\end{center}
\end{merksatz}
